\bookchapter{第十一章：攝影與圖片管理}

\chapterintro{說句實話，大多數人用手機拍照，拍出來的效果遠遠不如手機本身的潛力。不是你技術不行，而是很多小設定和小習慣你沒注意到。這一章不講虛的，就說那些你能馬上用上的東西。}

\booksection{一、你是不是也犯過這些錯？}

先說說最常見的幾個問題，你看看中了幾條：

\begin{itemize}
  \item \textbf{鏡頭是髒的}：這個真的是最普遍的問題。掏手機的時候指紋按在鏡頭上，拍出來霧濛濛的，還以為手機壞了。養成習慣，拍照前用衣角擦一下鏡頭。
  \item \textbf{從來不調曝光}：預設模式下，點一下螢幕是對焦，但如果你在點完之後上下滑動，是可以調節亮度的。很多人拍的照片要麼太暗要麼太亮，就是不知道這個操作。
  \item \textbf{閃光燈永遠開著}：手機LED閃光燈拍出來的效果真的很醜，光線又硬又冷，把人拍得像鬼。除非真的伸手不見五指，否則關掉它。暗光環境下，用夜景模式遠比閃光燈效果好。
  \item \textbf{放大靠兩根手指拉伸}：在螢幕上雙指縮放是數碼變焦，本質就是裁切畫面然後放大，畫質會明顯下降。如果想拍遠處的，直接用取景框上的變焦按鈕（0.5x/1x/2x/3x）。其中 0.5x（超廣角）和 1x（主攝）是原生光學鏡頭，不損失畫質；2x/3x 則多為"光學品質變焦"——透過高畫素感測器裁切實現，畫質接近光學變焦但並非獨立鏡頭。超過 3x 的變焦通常就是混合變焦或數碼變焦了，畫質會有明顯下降。
\end{itemize}

\booksection{二、幾個立刻提升出片率的小技巧}

這些不需要什麼理論知識，用一次就能看到效果：

\begin{itemize}
  \item \textbf{三分法構圖+網格線}：在相機設定裡開啟"網格線"（也叫參考線）。拍照時把你想突出的東西放在四條線的交叉點上，畫面立刻就會"對"很多。這是攝影裡最簡單也最有效的一個規則，沒什麼玄學，就是符合人眼的視覺習慣。
  \item \textbf{找一找"引導線"}：拍風景或者街景的時候，留意一下畫面裡有沒有自然的線條——馬路、欄杆、牆壁、一行樹——讓這些線條從畫面邊緣延伸到主體。這個技巧用好了，路人看你拍的照片會以為你學過攝影。
  \item \textbf{別老站著拍}：大部分人拍照永遠是自己眼睛的高度，所以拍出來的視角千篇一律。試試蹲下來從低角度拍，或者找個高處往下拍，同一個地方出來的感覺完全不一樣。拍小孩和寵物尤其需要蹲下來，平視才能拍出自然的感覺。
  \item \textbf{光線是第一要素}：中午大太陽下拍人像，臉上的陰影又硬又醜。最好的光線是清晨和傍晚的"黃金時段"，陽光是斜的、顏色偏暖，拍什麼都好看。陰天的光線雖然不出彩，但勝在均勻柔和，拍人像反而比大太陽下容易出片。
  \item \textbf{夜景模式使用要領}：開啟夜景模式之後，手機會自動連拍多張然後合成。這時候最關鍵的是——手絕對不能抖。能靠牆就靠牆，能放桌上就放桌上，實在不行就屏住呼吸。夜景模式需要你穩住2到4秒，抖了畫面就糊了。
\end{itemize}

\booksection{三、人像模式與專業模式}

\begin{itemize}
  \item \textbf{人像模式}：它能拍出背景模糊、主體清晰的效果（淺景深）。手機因為鏡頭物理尺寸小，真實虛化能力有限，所以靠的是AI演算法來實現——簡單說就是AI把你的人摳出來，然後模糊背景。用的時候保持1-2米距離效果最好，離太遠或太近演算法都容易出錯。還有，拍人像的時候，注意背景裡不要有太複雜的線條，否則AI扣圖會在頭髮邊緣出現"翻車"。
  \item \textbf{專業模式}：如果你想要完全掌控畫面效果，可以試試專業模式。幾個關鍵引數要知道：ISO越低畫質越乾淨，但畫面會暗；快門越慢進光越多，但手抖會糊；白平衡管色彩的冷暖，調不好照片會偏藍或偏黃。一般人不需要常年用專業模式，但瞭解這幾個引數至少能讓你知道相機在幹什麼。
\end{itemize}

\booksection{四、後期修圖：不要過度，恰到好處就行}

很多人對修圖有誤解，覺得修了就不真實了。其實相機拍出來的照片本來就和肉眼看到的不一樣，適當調整是為了還原你當時看到的畫面，而不是造假。

\begin{itemize}
  \item \textbf{Snapseed}：谷歌出的免費修圖工具，功能專業但沒有廣告。我日常修圖的流程是：先用"調整圖片"拉一拉亮度和對比度 用"突出細節"銳化一下 用"白平衡"糾正偏色 最後"裁剪"成想要的構圖。Snapseed的"區域性"功能很好用，可以只調畫面中某一個區域，不用全圖一起動。
  \item \textbf{醒圖/美圖秀秀}：這兩個適合修人像。但我覺得有個原則：美顏濾鏡拉到30\%以內就好，超過這個數就假了，皮膚像塑膠一樣。自然比完美好看。
  \item \textbf{Lightroom手機版}：如果你對色彩有追求，這個是最專業的。它有曲線調整、HSL分色調節這些高階功能。不過對大多數人來說有點殺雞用牛刀。
  \item \textbf{修圖工作流建議}：我自己的習慣是先拍好原片，回家後在電腦螢幕或iPad上慢慢修，手機螢幕太小，色彩也不太準。如果必須在手機上修，記得把螢幕亮度調到中等偏上，別在省電模式下修圖。
\end{itemize}

\booksection{五、照片太多怎麼辦？整理策略}

手機相簿裡堆了幾千張照片，想找一張特定的翻半天——這個體驗大家都有。幾個實用的整理方法：

\begin{itemize}
  \item \textbf{善用搜尋}：現在iPhone和大部分安卓手機的相簿都支援智慧搜尋。你可以搜"貓""海邊""身份證""花"這些關鍵詞，系統會用AI識別畫面內容幫你找出來。還可以搜日期和地點，比如"去年五月""北京"。這個功能知道的人不多，但真的很好用。
  \item \textbf{定期清理垃圾}：截圖、重複的照片、拍糊了的廢片，這些才是佔空間的元兇。iOS相簿有"重複專案"自動合併功能；安卓可以去檔案管理裡搜"screenshots"資料夾，批次刪不需要的截圖。我每個月花五分鐘清理，就能騰出好幾個G。
  \item \textbf{建立相簿，別全靠預設排序}：系統預設按時間排序沒錯，但對於"孩子的成長照""旅行照片""重要證件"這種跨時間段的內容，你應該手動建相簿歸類。蘋果的"共享相簿"還能跟家人一起往裡面加照片，一個人建，全家都能看。
  \item \textbf{備份是底線}：照片是最不可替代的資料——檔案丟了可以重建，聊天記錄沒了也就沒了，但孩子的照片、老人的影片丟了就真的沒了。一定要開雲備份。iCloud照片最省心，安卓這邊各品牌都有自己的雲服務，或者用百度網盤/阿里雲盤自動備份。原則是：至少有一份雲備份，別把雞蛋放在同一個籃子裡。
\end{itemize}

\booksection{六、截圖與錄屏}

\begin{itemize}
  \item \textbf{截圖}：同時按電源鍵+音量加鍵，所有手機通用的方法。很多安卓手機還支援三指下滑截圖、指關節雙擊截圖。長截圖（滾動截圖）在安卓上很成熟——截完一張後點"滾動"按鈕，手機會自動往下翻頁並拼接成一張長圖，適合截聊天記錄和網頁。
  \item \textbf{錄屏}：從控制中心新增"螢幕錄製"，點一下開始錄，再點停止。錄之前長按錄製按鈕可以選擇是否錄製麥克風聲音——你要做教程解說就開麥，只是錄操作畫面就關掉。錄完的影片會自動存到相簿裡。
\end{itemize}

\begin{tcbpractice}{練一練：用手機完成一次完整的拍照+修圖+整理}
\begin{enumerate}
  \item 開啟相機設定裡的"網格線"，找一個有自然光的地方拍三張照片（注意擦鏡頭、調曝光）。
  \item 選最好的一張，用Snapseed做基礎調整（亮度、對比度、裁剪）。
  \item 給這張照片建一個專屬相簿，命名為"我的作品"。
  \item 檢查一下雲備份有沒有開啟，確保照片已經在雲端存了一份。
\end{enumerate}

\end{tcbpractice}
